\documentclass[11pt,a4paper]{article}
\usepackage[T1]{fontenc}
\usepackage[utf8]{inputenc}
\usepackage[main=italian,provide=*]{babel}
\usepackage{amsmath,amssymb}
\usepackage{geometry}
\geometry{margin=2.3cm,top=2.5cm,bottom=2.7cm}
\usepackage{booktabs}
\usepackage[table]{xcolor}
\usepackage{tcolorbox}
\tcbuselibrary{skins,breakable}
\usepackage{titlesec}
\usepackage{enumitem}
\usepackage{needspace}
\usepackage{fancyhdr}
\usepackage{hyperref}
\hypersetup{colorlinks=true,linkcolor=Sepia,citecolor=Sepia,urlcolor=Sepia}

% --- palette ---
\definecolor{inkblue}{HTML}{1B3A5C}
\definecolor{lightblue}{HTML}{EAF1F8}
\definecolor{accent}{HTML}{B0562A}
\definecolor{softgray}{HTML}{6B6B6B}
\definecolor{okgreen}{HTML}{2E6B3E}
\definecolor{okgreenbg}{HTML}{EAF5EC}
\definecolor{advisorpurple}{HTML}{5B3A7A}
\definecolor{advisorbg}{HTML}{F1ECF6}

% --- titoli di sezione ---
\titleformat{\section}
  {\normalfont\Large\bfseries\color{inkblue}}
  {\thesection}{0.6em}{}[{\color{inkblue!40}\titlerule}]
\titlespacing{\section}{0pt}{0.7em}{0.4em}
\titleformat{\subsection}
  {\normalfont\large\bfseries\color{accent}}
  {}{0em}{}
\titlespacing{\subsection}{0pt}{0.5em}{0.2em}

\pagestyle{fancy}
\fancyhf{}
\renewcommand{\headrulewidth}{0.4pt}
\fancyhead[L]{\small\color{softgray}Modello di rumore --- Passo 1, Parte 2}
\fancyhead[R]{\small\color{softgray}\thepage}
\fancyfoot[C]{\small\color{softgray}Tesi triennale, Samuele --- Relatore: Prof.\ Alessandro Chiesa}

% --- box riutilizzabili ---
\newtcolorbox{keyeq}[1][]{
  colback=lightblue, colframe=inkblue!70, boxrule=0.6pt,
  arc=2pt, left=8pt, right=8pt, top=3pt, bottom=3pt, #1
}
\newtcolorbox{panelbox}[1]{
  colback=white, colframe=softgray!50, boxrule=0.5pt,
  arc=2pt, left=7pt, right=7pt, top=2pt, bottom=2pt,
  fonttitle=\bfseries\color{inkblue}, title={#1}
}
\newtcolorbox{okbox}[1]{
  colback=okgreenbg, colframe=okgreen!60, boxrule=0.6pt,
  arc=2pt, left=7pt, right=7pt, top=2pt, bottom=2pt,
  fonttitle=\bfseries\color{okgreen}, title={#1}
}

\title{\vspace{-1.5em}\color{inkblue}Modello di rumore del dimero --- Passo 1\\[0.2em]
\Large\normalfont\color{softgray}Calibrazione reale e verifica del limite di rumore nullo}
\author{Note di lavoro --- Tesi triennale, Samuele\\ Relatore: Prof.\ Alessandro Chiesa}
\date{}

\begin{document}
\sloppy
\maketitle
\thispagestyle{fancy}
\vspace{-2.2em}

File: \texttt{noise\_model\_dimero.py}, \texttt{validate\_noise\_model\_dimero.py}.
Teoria di riferimento: \texttt{teoria\_modello\_rumore.tex}.

\section*{Metodologia (nessuna scelta nuova)}
\addcontentsline{toc}{section}{Metodologia}

\begin{itemize}[leftmargin=1.3em,itemsep=0.25em]
\item \textbf{Due soli canali}, confermati dal relatore il 02/09/2026: errore
di gate (depolarizzante, 1 e 2 qubit) ed errore di lettura (readout
simmetrico). $T_1$/$T_2$ esclusi.
\item \textbf{Conversione} $\varepsilon\to\lambda$ derivata in
\texttt{teoria\_modello\_rumore.tex} a partire dalla fedeltà media del
canale depolarizzante:
\begin{keyeq}
\centering
$\lambda_{1q}=2\varepsilon_{1q}, \qquad \lambda_{2q}=\dfrac43\varepsilon_{2q}$
\end{keyeq}
\item \textbf{Transpilazione per blocco}: il circuito viene transpilato una
volta e il blocco compilato riusato, non ritranspilato per intero --- vedi
la trappola isolata nella sessione precedente (\texttt{log\_decisioni.md}).
Per questo documento, che riguarda solo il VQE (nessun passo di Trotter
ripetuto), la distinzione non è ancora rilevante: torna cruciale al Passo 3.
\end{itemize}

\section*{Calibrazione scelta}
\addcontentsline{toc}{section}{Calibrazione scelta}

\textbf{Chip:} \texttt{ibm\_torino} (IBM Heron r1, 133 qubit). \textbf{Fonte:}
J.~P.~T.~Stenger, G.~Bazargan, N.~T.~Bronn, D.~Gunlycke, \emph{“Method for
simulating open-system dynamics using mid-circuit measurements on a quantum
computer”}, arXiv:2504.15187 (2025), Appendice B --- mediane pubblicate.

\begin{center}
\renewcommand{\arraystretch}{1.25}
\begin{tabular}{@{}lcc@{}}
\toprule
\textbf{Grandezza} & \textbf{Valore (mediana)} & \textbf{Origine} \\
\midrule
$\varepsilon_{1q}$ & $2.9\times10^{-4}$ & gate $\sqrt X$ \\
$\varepsilon_{2q}$ & $3.8\times10^{-3}$ & gate CZ \\
$p_\text{readout}$ & $2.3\times10^{-2}$ & errore di lettura \\
\midrule
$\lambda_{1q}=2\varepsilon_{1q}$ & $5.80\times10^{-4}$ & derivato \\
$\lambda_{2q}=\tfrac43\varepsilon_{2q}$ & $5.067\times10^{-3}$ & derivato \\
\bottomrule
\end{tabular}
\end{center}

\noindent \textbf{Nota di modellazione:} il gate nativo a 2 qubit di
\texttt{ibm\_torino} è CZ, mentre i nostri circuiti sono scritti in base
$\{R_z,\sqrt X,X,\text{CX}\}$. Si usa $\varepsilon_{2q}$ come valore
rappresentativo dell'ordine di grandezza dell'errore a 2 qubit --- coerente
con l'indicazione del relatore (``usa dei valori tipici... di qualche chip
di riferimento''), non una replica esatta dell'hardware. Lo scan del Passo 5
copre comunque un intorno di questo valore, non solo il punto singolo.

\section*{Verifica: limite di rumore nullo}
\addcontentsline{toc}{section}{Verifica: limite di rumore nullo}

Controllo obbligatorio prima di usare il modello su un circuito completo
(dichiarato in \texttt{teoria\_modello\_rumore.tex}): con
$\varepsilon_{1q}=\varepsilon_{2q}=p_\text{readout}=0$, il circuito VQE
(ansatz PMA-2q$\cdot$3, parametri da \texttt{ground\_state\_test2.npz}) su
\texttt{AerSimulator(method="density\_matrix")} deve riprodurre esattamente
$E_\text{VQE}$ registrato in Parte 1.

\begin{okbox}{Limite di rumore nullo verificato}
\begin{center}
\renewcommand{\arraystretch}{1.25}
\begin{tabular}{@{}lc@{}}
\toprule
\textbf{Quantità} & \textbf{Valore}\\
\midrule
$E_\text{VQE}$ (registrato, Parte 1) & $-3.57321145$\\
$E_\text{VQE}$ (density\_matrix, rumore nullo) & $-3.57321145$\\
$|\Delta E|$ & $3.46\times10^{-9}$\\
CNOT nel circuito transpilato & $2$ (invariato)\\
\bottomrule
\end{tabular}
\end{center}
\end{okbox}

\section*{Effetto del rumore di riferimento}
\addcontentsline{toc}{section}{Effetto del rumore di riferimento}

Con i valori di calibrazione di \texttt{ibm\_torino} (tabella sopra),
l'energia VQE stimata sul simulatore rumoroso si allontana da quella ideale
di una quantità fisicamente sensata:

\begin{center}
\renewcommand{\arraystretch}{1.2}
\begin{tabular}{@{}lc@{}}
\toprule
\textbf{Quantità} & \textbf{Valore}\\
\midrule
$E_\text{VQE}$ (rumore di riferimento) & $-3.50164605$\\
scostamento da $E_\text{VQE}$ ideale & $+0.071565$\\
\bottomrule
\end{tabular}
\end{center}

\section*{Scan preliminare su $\varepsilon_{2q}$}
\addcontentsline{toc}{section}{Scan preliminare}

A titolo di controllo di sanità (non ancora lo scan completo del Passo 5,
che coprirà anche $\varepsilon_{1q}$ e $p_\text{readout}$): a parità di
$\varepsilon_{1q}=2.9\times10^{-4}$ e $p_\text{readout}=2.3\times10^{-2}$,
lo scostamento da $E_\text{VQE}$ ideale cresce in modo monotono con
$\varepsilon_{2q}$, come atteso.

\begin{center}
\renewcommand{\arraystretch}{1.2}
\begin{tabular}{@{}cccc@{}}
\toprule
$\varepsilon_{2q}$ & $E_\text{VQE}$ rumoroso & scostamento \\
\midrule
$0$ & $-3.53740090$ & $+0.035811$ \\
$1\times10^{-3}$ & $-3.52797412$ & $+0.045237$ \\
$3.8\times10^{-3}$ (riferimento) & $-3.50164605$ & $+0.071565$ \\
$8\times10^{-3}$ & $-3.46233882$ & $+0.110873$ \\
$1.5\times10^{-2}$ & $-3.39731982$ & $+0.175892$ \\
\bottomrule
\end{tabular}
\end{center}

\noindent Anche a $\varepsilon_{2q}=0$ lo scostamento resta non nullo
($+0.036$): è il contributo di $\varepsilon_{1q}$ e $p_\text{readout}$, non
ancora annullati in questa riga della tabella.

\needspace{6\baselineskip}
\section*{Prossimo passo}
\addcontentsline{toc}{section}{Prossimo passo}

Colmare la lacuna già segnalata: implementare \texttt{ansatz\_params} nel
modulo dei correlatori del dimero (per ora dispone solo di
\texttt{prepare\_state}), così da poter agganciare questo stesso
\texttt{NoiseModel} anche lì. Poi, Passo 2 della pipeline: VQE con termine
DM sotto rumore come stadio a sé (energia e fedeltà rumorose confrontate
con l'esatto), separato dal suo uso come preparazione per i passi
successivi.

\end{document}
